\subsection{Results}
\label{sec:results-P}

\subsubsection{Results reproducing original paper}

The reproduced linear evaluation and zero-shot results are reported in Table~\ref{tab:t2ab} (appendix). The DINOv2 row is the cleanest comparison (both checkpoints from the original paper) and reproduces the central claim: trained registers improve all three probes. OpenCLIP with test-time registers shows no Top-1 change and small dense-prediction gains, consistent with \citet{jiang2025vision}; the $-0.67$~pt zero-shot drop (one misclassified image on 150) is within subset variance. DeiT-III with injected \emph{untrained} register-token parameters regresses $-0.67$~pt on Top-1 but still improves mIoU and RMSE, suggesting that even random register slots act as a partial sink for high-norm outliers -- evidence in favour of the paper's mechanistic interpretation that registers' role is to absorb artifact tokens rather than to carry trained information.

\subsubsubsection{Figure 8 -- Attention maps and linear evaluation vs.\ number of registers}

\citet{darcet2024vision} sweep $N\in\{0,1,2,4,8,16\}$ register tokens by retraining DINOv2 from scratch for each value -- a regime out of our compute budget. We instead compare three retraining-free approximations defined in Section~\ref{sec:method-P}. Qualitative attention maps for all three are stacked in Figure~\ref{fig:f8-top} (appendix); the corresponding downstream metrics are in Figure~\ref{fig:f8-bot} and Table~\ref{tab:f8} (appendix). The truncate/cycle proxy collapses at $N=0$ (Top-1 $70.67\%$, since DINOv2-reg4 expects exactly four registers) and plateaus from $N=1$ onwards. The Jiang and independent variants on the unmodified baseline both jump from $94.00\%$ Top-1 at $N=0$ to $95.33\%$ at $N=1$, plateau through $N=8$, and then collapse at $N=16$ ($92.00\%$ for Jiang, $92.67\%$ for the independent variant). 

\subsubsection{Results beyond original paper}

\subsubsubsection{Independent variant of the Jiang test-time register intervention}

A close reading of Jiang's code (\texttt{shared/hook\_fn.py}, function \texttt{activate\_on\_registers}) reveals that for each register neuron the top-1 outlier value is duplicated identically across all $N$ test-time registers. The collapse we observe at $N=16$ might therefore be caused by the CLS token attending to 16 near-identical broadcast nodes and diluting its global readout. To test this we introduced a controlled variant in which, for each neuron, the top-$N$ distinct outlier patches are picked and their signed values are written onto $N$ different registers (so register $i$ receives the $(i{+}1)$-th outlier). Patch positions are zeroed exactly as in Jiang's pipeline. The numerical results are reported alongside the Jiang re-implementation in Table~\ref{tab:f8} (appendix); the qualitative attention maps are in Figure~\ref{fig:f8-top} row~3. \textbf{The hypothesis is falsified}: the $N=16$ collapse persists ($92.67\%$ vs.\ $92.00\%$ for Jiang), the Top-1 profile is otherwise identical to Jiang's, and diagnostic measurements confirm that the patch-token norms after the intervention are clean at all $N\geq 1$ (max patch norm $\sim 51$, no patch above $60$). Since neither the writing policy nor residual artifacts on the patches explain the collapse, the failure mode must lie in the downstream transformer blocks (18--23 of DINOv2 ViT-L/14) that follow the intervention point and were never trained to process so many simultaneously active broadcast tokens. This is a falsifiable finding that does not appear in either source paper.
